\documentclass[12pt]{article}

\usepackage[utf8]{inputenc}
\usepackage[T1]{fontenc}
\usepackage{lmodern}
\usepackage{csquotes}
\usepackage{amsmath, amssymb, amsthm, bm}
\usepackage{setspace}
\usepackage{geometry}
\usepackage{xurl}
\usepackage{hyperref}

\geometry{margin=1in}
\setstretch{1.2}

\title{Visions and Boundaries: Swedenborg, Kant, and the Modern Crisis of Mind}
\author{Flyxion}
\date{\today}

\begin{document}

\maketitle

\begin{abstract}
This extended essay explores the historical interplay between the visionary 
world of Emanuel Swedenborg (1688--1772) and the critical philosophy of 
Immanuel Kant (1724--1804), examining how their confrontation shaped modern 
theories of mind, perception, and the limits of reason. Swedenborg's transition 
from Enlightenment scientist to mystic seer produced a theology of 
correspondences, apocalyptic revelations, and a detailed phenomenology of 
visions. Kant's response---first fascinated, then recoiling through the ironic 
\emph{Dreams of a Spirit-Seer}---served as a catalyst for the epistemological 
boundaries later formalized in the \emph{Critique of Pure Reason}. The essay 
connects these eighteenth-century debates to modern research on aphantasia, 
hyperphantasia, hallucinations in large language models, and the global 
metacrisis of sense-making. The result is a unified historical and conceptual 
analysis of how the tension between unbounded imagination and disciplined 
reason continues to shape cognitive, philosophical, and technological 
landscapes today.
\end{abstract}

\section{Swedenborg: The Scientist Who Became a Seer}

Emanuel Swedenborg's life divides sharply into two phases: the empirical 
investigator and the visionary prophet. During his early decades, Swedenborg 
pursued work recognizable as modern scientific research. He studied under 
Christopher Polhem, Sweden's leading engineer, and contributed to designs for 
canals, mining pumps, and mechanical innovations. His scientific publications 
included treatises on metallurgy, mineralogy, cosmology, and the circulation 
of blood.

In these early works, Swedenborg sought what he called ``the harmonies of the 
universe''---a structural order underlying the diversity of natural 
phenomena. His anatomical speculations attempted to locate the ``seat of the 
soul'' in the pineal gland and cortical fibers, anticipating certain themes 
in later neuroscience.

\subsection{The Crisis of 1744--1745}

The turning point came during travels in England and the Netherlands. In 1744 
Swedenborg began experiencing vivid dreams and intense somatic disturbances. 
His \emph{Journal of Dreams}, discovered posthumously, documents encounters 
with radiant beings, terrifying symbolic figures, and overwhelming ``waves of 
heat'' he interpreted as divine contact.

By 1745, he believed himself commissioned by Christ to reveal the true 
spiritual meaning of Scripture. He described the physical and spiritual worlds 
as interacting planes, the latter populated by angels, spirits, and the 
departed.

\subsection{The Spiritual Diary and the Afterlife}

Swedenborg’s \emph{Arcana Coelestia} (1749--1756), an eight-volume exegesis of 
Genesis and Exodus, inaugurated his visionary theology. It presents a system 
of \emph{correspondences}: every material event expresses a spiritual analogue. 
This interpretive method extends to nature, psychology, and cosmology.

His 1758 work \emph{Heaven and Hell} provides a detailed ethnography of the 
afterlife. Angels appear in communities structured by love and wisdom; spirits 
travel through regions shaped by their inner states. Space and time in the 
spiritual world, he claimed, are ``states of mind'' rather than physical 
dimensions.

\subsection{The Apocalyptic Judgment of 1757}

In one of his most dramatic claims, Swedenborg announced that the Last Judgment 
occurred in the spiritual world in 1757. This ``judgment'' did not manifest in 
earthly catastrophe but in the reordering of spiritual societies. The old 
Christian forms, he insisted, had spiritually collapsed; a new era of internal 
religion had begun.

\section{Kant Before the Critical Turn}

Immanuel Kant, born 36 years after Swedenborg, was shaped by a very different 
temperament. His early work spanned natural philosophy, cosmology, and 
anthropology. He wrote on planetary formation, earthquakes, the physics of 
wind, and the structure of the heavens. His 1755 \emph{Universal Natural 
History and Theory of the Heavens} proposed nebular hypotheses anticipating 
later astrophysics.

\subsection{Kant’s Encounter with Swedenborg}

In the 1750s, reports of Swedenborg's clairvoyant feats---especially the 
alleged prediction of the Stockholm fire while he was in Gothenburg---reached 
Kant. Curious, he obtained volumes of the \emph{Arcana Coelestia}. The 
purchase was expensive; the reading was unsettling. Kant found Swedenborg’s 
visions alternately intriguing, absurd, and disturbing.

Letters from the period reveal a philosopher wrestling with the boundaries of 
reason. Swedenborg threatened Kant’s commitment to Enlightenment rationality: 
if such visions were true, then the world was stranger than reason could admit.

\subsection{\emph{Dreams of a Spirit-Seer}}

Kant’s 1766 treatise, \emph{Träume eines Geistersehers}, is a hybrid text. 
Presented as satire, it contains a series of deep metaphysical reflections. 
Kant ridicules Swedenborg’s ``spirit geography'' yet simultaneously explores 
the possibility of a non-sensible realm accessible only through pure reason.

The key move is Kant’s separation of:
\begin{enumerate}
    \item the world as we experience it (the \emph{phenomenal} realm), and  
    \item the world as it exists independently of our cognitive structures (the \emph{noumenal} realm).
\end{enumerate}

This distinction becomes foundational in the \emph{Critique of Pure Reason} 
(1781), where Kant argues that space, time, and causality are not things in 
themselves but forms of intuition and categories of understanding. Swedenborg, 
for Kant, exemplified the peril of overstepping these boundaries.

\section{Swedenborg and Kant as Cognitive Extremes}

The Swedenborg--Kant relationship is illuminating when viewed through modern 
cognitive science. Psychological research identifies a spectrum of mental 
imagery from \emph{aphantasia} (lack of imagery) to \emph{hyperphantasia} 
(extremely vivid imagery).

\subsection{Kant as Conceptual Thinker}

Kant’s biography suggests a weak visual imagination. His thought was 
structural, architectural, and conceptual. He seldom used imagery, instead 
favoring abstraction and rigorous logical organization.

His need to impose \emph{limits} on cognition parallels the cognitive profile 
of aphantasia: ideas emerge not from images but from relations, constraints, 
and boundaries.

\subsection{Swedenborg as Visionary Imagination}

Swedenborg’s visions align with hyperphantasia:  
vivid imagery, multisensory hallucination, symbolic correspondences, and 
dream-like continuity.

His visionary system behaves like a generative model unconstrained by 
empirical grounding. Once a symbolic pattern appears, it cascades into 
self-reinforcing elaborations.

\section{Hallucination, Generative Models, and Kantian Alignment}

The tension between Swedenborg and Kant parallels the contemporary problem of 
LLM hallucinations. Large language models generate fluent but sometimes false 
outputs when internal statistical associations diverge from grounded truth.

\subsection{Swedenborg as Generative System}

Swedenborg’s imagination operated autopoietically. His correspondences grew 
in complexity without external constraint. In modern terms, his visions 
resembled autoregressive sampling from a richly parameterized internal model.

Just as LLMs can hallucinate plausible but ungrounded narratives, Swedenborg 
constructed an internally consistent symbolic universe detached from empirical 
verification.

\subsection{Kant as the First Alignment Architect}

Kant’s critical philosophy functions as a cognitive alignment system. It 
restricts the domain of valid inference to the conditions of possible 
experience.

Kant’s principles---unity, causality, substance, necessity---can be seen as 
analogues of constraints that prevent generative models from drifting into 
fantasy. His philosophy is an early attempt to balance creativity with 
epistemic safety.

\section{The Metacrisis: Vision Without Ground, Structure Without Flexibility}

The metacrisis refers to the interwoven crises of ecology, politics, 
technology, epistemology, and institutional legitimacy. Both Swedenborg and 
Kant illuminate aspects of this crisis.

\subsection{The Proliferation of Unbounded Vision}

The modern information ecosystem resembles a Swedenborgian cosmos:  
correspondences proliferate without constraint; symbolic systems flourish 
without grounding; conspiracies and ideologies form self-justifying spirals.

This is the danger of hyperphantasic society: imaginative excess without 
verification.

\subsection{The Rigidity of Systems}

Conversely, Kant’s insistence on rigid categories mirrors the technocratic 
world’s tendency toward over-systematization. Bureaucratic rationality, 
algorithmic governance, and ``solutionism'' reduce reality to narrow models.

Systems become brittle; they fail to adapt to novelty.

\subsection{Balancing Revelation and Reason}

The future requires neither Swedenborg's unbounded imagination nor Kant’s 
absolute constraints alone. It requires a dynamic synthesis:  
\begin{quote}
Imagination disciplined by structure,  
Structure permeated by imagination.
\end{quote}

This interplay is central to addressing the metacrisis.

\section{Conclusion}

Swedenborg and Kant represent cognitive archetypes: the visionary who sees too 
much and the architect who builds the walls that protect us from seeing too 
much. Their confrontation prefigures modern tensions between generative 
creativity and aligned reasoning, between symbolic overflow and epistemic 
boundaries.

In a time of global cognitive breakdown, their lives offer lessons in how 
vision and structure must be integrated. Swedenborg teaches the power and 
danger of imagination; Kant teaches the necessity and limit of reason. Our era demands both.

\newpage
\begin{thebibliography}{99}

\bibitem{swedenborg_heavenhell}
Swedenborg, Emanuel. \emph{Heaven and Hell}. 1758.

\bibitem{swedenborg_arcana}
Swedenborg, Emanuel. \emph{Arcana Coelestia}. 1749--1756.

\bibitem{swedenborg_journal}
Swedenborg, Emanuel. \emph{Journal of Dreams}. Manuscript.

\bibitem{jonsson_swedenborgbio}
Jonsson, Inge. \emph{Emanuel Swedenborg}. 2013.

\bibitem{larsen_swedenborg}
Larsen, Timothy. \emph{A People of One Book}. 2011.

\bibitem{kuehn_kant}
Kuehn, Manfred. \emph{Kant: A Biography}. Cambridge University Press, 2001.

\bibitem{cassirer_kant}
Cassirer, Ernst. \emph{Kant's Life and Thought}. 1981.

\bibitem{kant_dreams}
Kant, Immanuel. \emph{Dreams of a Spirit-Seer}. 1766.

\bibitem{kant_cpr}
Kant, Immanuel. \emph{Critique of Pure Reason}. 1781.

\bibitem{allison_freedom}
Allison, Henry. \emph{Kant's Theory of Freedom}. Cambridge University Press, 1990.

\bibitem{zeman_aphantasia}
Zeman, Adam et al. ``Lives Without Imagery: Congenital Aphantasia.'' \emph{Cortex}, 2015.

\bibitem{zeman_hyperphantasia}
Zeman, Adam et al. ``The Psychology of Hyperphantasia.'' 2020.

\bibitem{llm_hallucination1}
Ji, Ziwei et al. ``Hallucination in Natural Language Generation.'' \emph{ACM Computing Surveys}, 2023.

\bibitem{llm_hallucination2}
Rawte, Vinay et al. ``Survey of Hallucination in LLMs.'' \emph{arXiv}, 2023.

\bibitem{xu_interpretability}
Xu, Diyi et al. ``Understanding Model Hallucination.'' \emph{NeurIPS}, 2024.

\bibitem{han_psychopolitics}
Han, Byung-Chul. \emph{Psychopolitics}. Verso, 2017.

\bibitem{haraway_trouble}
Haraway, Donna. \emph{Staying with the Trouble}. Duke University Press, 2016.

\bibitem{tarnas}
Tarnas, Richard. \emph{The Passion of the Western Mind}. 1991.

\end{thebibliography}

\end{document}